\chapter{Patent Landscape}
\label{ch:patents}

Patent publications expose proposed measurement architectures that can be hard to reconstruct from product brochures.
Their technical value lies in the signal being observed, the calibration and inference steps, and the conditions under which an output is identifiable.
A disclosure, a granted claim, a commercial implementation and an accuracy experiment are four different kinds of evidence.
This chapter compares selected disclosures; \cref{app:patents} provides a wider catalog of document identifiers.

Read an abstract as an overview, a description as a set of disclosed embodiments, and each claim as its own combination of limitations.
An example camera count, exposure time or receiver arrangement in the specification is not automatically required by every claim.
Conversely, a broad title such as ``spin measurement'' does not mean that a patent covers every method producing that output.
Company groupings below aid navigation; they do not replace assignment records or establish which product implements a particular embodiment.

\section{TrackMan A/S (Interactive Sports Games A/S) --- Fredrik Tuxen}
\label{sec:patenttrackman}

TrackMan's virtual-marking notice contains 47 distinct US patent numbers at the October 1, 2026 check~\cite{trackmanpatents}.
That dated list is a route into its disclosures, not an exhaustive search of applications, foreign publications or the company's ownership history.

\begin{description}
\item[Spin Frequency and Trajectory-Based Axis.]
\patent{US8845442B2}.
The disclosure describes radar spectral spacings associated with spin
and a separate trajectory-based approach to spin axis~\cite{us8845442}.
Claim~1 concerns estimating spin frequency from a frequency distance
in the modulated radar signal; it should not be paraphrased as requiring
all the axis-inversion steps in the description. Spectral structure
still needs a defensible mapping from the observed periodicity to ball
rotation. Trajectory-based axis inference also depends on the aerodynamic
model and the information present in the observed flight segment
(\cref{sec:radarspin,eq:liftextract}).
\item[Target-Referenced Measurements.]
\patent{US8085188B2} / \patent{US9857459B2} /
\patent{US10473778B2}.
These documents concern relating launch or trajectory information to a
target designated in an image. Sensor and image coordinates must be
registered before a target-relative angle has a physical meaning. Their
publication history does not by itself establish clearance for a
particular target-selection interface.
\item[A Radar View of Spin Axis.]
\patent{US10850179B2}.
Claim~1 separates the Doppler spectrum into components, determines their
angular positions using a non-collinear receiver arrangement and identifies
a \emph{projection} of the spin axis onto a plane perpendicular to the
radar's line of sight~\cite{us10850179}. That projection is not, by itself,
an unrestricted three-dimensional axis measurement. Motion compensation,
receiver geometry and the distinction between a projected axis and a full
axis belong in any explanation of the method. Related documents include
\patent{US11446546B2} and \patent{US11938375B2}; their claims must be read
separately.
\item[Tracking, Calibration and Impact.]
\patent{US10989791B2} and \patent{US11828867B2} concern radar/imager
tracking; \patent{US11619708B2} and \patent{US12517218B2} address
inter-sensor calibration; \patent{US10953303B2} describes estimating
club impact information from image features. These are distinct inference
problems, even when a product presents their outputs together. The
similarly titled radar/image documents \patent{US10596416B2},
\patent{US11697046B2} and \patent{US12128275B2} identify Topgolf Sweden AB,
not TrackMan, in their publication records.
\item[Display and Attribution.]
Trajectory illustration appears in \patent{US10315093B2}.
Associating tracked balls with launch locations is described in \patent{US11086005B2}. A convincing video overlay and a correct association with a
player are separate requirements from accurate trajectory measurement.
\end{description}

\section{Foresight Sports, WAWGD and Wintriss Disclosures}
\label{sec:wintriss}

The Wintriss/WAWGD documents provide examples of image-based launch measurement and integrated monitor operation.
Uneekor's product marking notice lists several of the same numbers, and the companies announced a license agreement in 2024~\cite{uneekorpatents,businesswireforesight}.
That evidence should not be expanded into an assertion that every Foresight or Uneekor device implements every disclosed feature.

\begin{description}
\item[Surface Images.]
\patent{US7292711B2} / \patent{US7324663B2}.
The disclosures describe recovering ball motion from images of surface
features, including dimples and non-precision marks~\cite{us7292711}.
US7292711B2 claim~1 includes compensation for sensor misalignment relative
to a calibrated orientation. US7324663B2 claim~1 describes a search over
rotations, scaling and correlation scores, together with coordinate and
alignment calculations. Here, ``markerless'' means that prescribed
fiducials need not be applied; it does not mean that the measurement needs
no resolvable texture, lighting control or calibration. Correlation
quality and rotation ambiguity must be evaluated on the actual images.
\item[Trigger and Placement.]
\patent{US7497780B2} / \patent{US7641565B2}.
The former integrates trigger-zone monitoring, capture, parameter
calculation and output. The latter detects ball placement using images
and physical attributes. Reliable triggering determines which frames
reach a spin estimator; it is part of the measurement chain, rather than
merely a display convenience.
\item[Club Motion.]
\patent{US8951138B2}.
This disclosure addresses club path and face orientation around impact.
Its inclusion supports discussion of club measurement methods, not a
blanket assertion that four cameras or reflective face dots are
universally protected, or that one patent identifies a particular
commercial club-data system.
\end{description}

\section{FlightScope / EDH --- Henri Johnson}
\label{sec:patentflightscope}

\begin{description}
\item[Phase-Modulated Spin Signals.]
\patent{US9868044B2}.
Claim~1 derives spin from the period of a demodulated radar component
whose portions correspond to apparent speeds above and below a nominal
speed~\cite{us9868044}. The disclosure's dielectric-lens explanation and
surface-feature examples should be distinguished from that claim's
signal-processing requirements. Symmetry can cause a feature pattern to
repeat more than once per revolution; a measured modulation frequency
therefore needs interpretation before it becomes a spin rate.
\item[Receiver-Pair Timing.]
\patent{US10775492B2}.
Claim~1 uses demodulated Doppler signals from receiver pairs in
perpendicular planes, their relative time delays and receiver separations
to calculate an axis relative to a receiver plane~\cite{us10775492}.
This supplies rotation-related observations without first fitting a
long curved trajectory. It does not establish unconditional indoor
accuracy: signal quality, delay estimation, geometry and the stated axis
convention remain part of the result. A ratio of delays also needs
careful treatment near a vanishing denominator and when resolving angle
quadrants; a bare one-argument arctangent is not a complete implementation.
\item[Camera/Radar Tracking.]
\patent{US10338209B2}.
The document describes calibrated simultaneous tracking and combining
camera and radar information~\cite{us10338209}. Claim~1 specifically
includes using radar data to interpolate a three-dimensional camera
trajectory when one sensor fails to supply usable measurements. Recovery
through such a gap is an inference step whose uncertainty must be tested,
not a second independent observation of the missing interval.
\end{description}

\section{Acushnet (Titleist) --- Optical Measurement Architectures}
\label{sec:patentacushnet}

These documents illustrate several ways to assemble an optical launch monitor~\cite{us5501463,us6500073,us6758759}.
They do not reduce to one fixed camera-and-marker recipe.

\begin{description}
\item[Marked Club and Ball.]
\patent{US5501463A}.
Claim~1 uses at least two cameras and repeated exposures of contrasting
areas on both the striking instrument and the object. Image patterns
support estimation of club path and orientation. The conversion from
marker positions to a club-face quantity requires a defined relationship
between those markers and the club geometry.
\item[Launch Conditions to Trajectory.]
\patent{US6500073B1}.
Claim~1 obtains positions and orientations from images at different times,
calculates launch conditions, then calculates a trajectory. Measured
initial conditions and the resulting modeled flight must be distinguished:
a predicted landing point is not an observed landing point.
\item[Separate Club and Ball Capture.]
\patent{US6758759B2}.
The first claim places club and ball cameras in vertically separated
planes and derives their motion data from repeated images. The minimum
in that claim is one camera for each task; describing the entire patent
as requiring two stereo pairs would confuse an embodiment with the claim.
\item[Integrated Optical Hardware.]
\patent{US7143639B2}.
Claim~1 specifies a housing layout with a club camera, at least two ball
cameras, illumination, triggering and a teeing aid, arranged for use by
both handednesses through rotation of the apparatus. The first claim does not require four cameras.
\item[Three-Dimensional Image Generation.]
\patent{US10668350B2}.
The disclosure discusses stereo and light-field approaches. Claim~1
specifically includes a light-field camera that captures light intensity
and ray direction, followed by reconstruction and kinematic processing.
Example frame rates or exposures in a disclosure are design proposals;
they do not demonstrate a product's achieved measurement error.
\end{description}

\section{United States Golf Association --- Aerodynamic Experiments}

The United States Golf Association's \patent{US6186002B1} belongs outside the Acushnet grouping.
It describes controlled launches and timed light-screen crossings for aerodynamic coefficient fitting (\cref{app:patents})~\cite{us6186002}.
Its experimental setup matters: an arbitrary short trajectory need not contain enough information to separate initial spin, wind and uncertain aerodynamic coefficients.

\section{Korean Ecosystem: Creatz/Uneekor and Golfzon}

Uneekor's marking notice connects specific products with listed patent numbers~\cite{uneekorpatents}.
The Creatz disclosures include \patent{US10247553B2}, which describes a start sensor for identifying ball placement and starting movement, and \patent{US12008770B2}, which concerns dimple-based spin estimation.
A product's markerless-spin label does not specify its entire algorithm or establish performance across every ball surface.

Creatz's \patent{US11191998B2} makes the distinction between observation and inference particularly clear.
Claim~1 uses speed- or angle-related quantities to estimate spin when image marks cannot supply it; when mark-based calculation is possible, it also uses such an estimate to correct that calculation.
This is more than a fallback for missing observations.
The reported spin can depend on both image evidence and a relationship inferred from other ball or club quantities.
Validation should distinguish these operating modes and test their errors separately; a plausible corrected value does not establish an independent measurement of rotation.

Golfzon's \patent{US12002222B2} provides a different example of this dependence.
Claim~1 obtains dynamic loft and backspin through a reference database using observed club/ball motion parameters, then calculates spin-axis and side-spin quantities from motion vectors and specified relationships.
These inferred outputs must not be used as independent evidence that the assumed impact relationships hold for a particular shot.

Golfzon's \patent{US9242158B2} describes beginning the displayed flight with a first set of ball information and updating it when a second set arrives~\cite{us9242158}.
This separates feedback latency from completion of the measurement pipeline.
Timing examples in a description are not measured latency guarantees.
Readers comparing monitors should ask both when an output appears and when the final parameters become available.

\section{Technical Research Map and Rights Boundaries}
\label{sec:fto}

The following map identifies engineering questions to investigate.
It replaces a technique-wide ``free to use'' classification: the age or expiration of one document cannot establish clearance for every implementation bearing the same technique name.
The USPTO describes a patent as a right to exclude, not permission to practice an invention.
Term cannot be computed indiscriminately as twenty years from any date labeled ``priority''; the applicable filing basis and other term conditions matter.\footnote{USPTO, \emph{Patent Essentials}, \url{https://www.uspto.gov/patents/basics/essentials}.} A design-specific rights review requires the relevant claims, jurisdictions and current records.
This technical catalog does not supply that determination.

\begin{table}[htbp]
\centering\small
\caption{Selected Measurement Methods and Their Validation Questions}
\label{tab:fto}
\begin{tabular}{@{}>{\raggedright\arraybackslash}p{3.5cm}>{\raggedright\arraybackslash}p{4.1cm}>{\raggedright\arraybackslash}p{6.6cm}@{}}
\toprule
\textbf{Method} & \textbf{Example Documents} & \textbf{Question to Test} \\
\midrule
Surface-image spin & US7292711B2; US7324663B2 &
  Are texture, exposure and rotation matching sufficient to resolve ambiguity? \\
Marked club/ball images & US5501463A; US6758759B2 &
  Are calibration and marker-to-face geometry known at the reported time? \\
Radar spin periodicity & US8845442B2; US9868044B2 &
  Which spectral or temporal repeat corresponds to one revolution? \\
Radar axis information & US10850179B2; US10775492B2 &
  Is the result a projection or a full axis, and in which coordinate frame? \\
Camera/radar fusion & US10338209B2; US10989791B2 &
  Are clocks, coordinates, association and missing-data uncertainty controlled? \\
Trajectory-informed spin & US12544624B2 &
  How strongly does the result depend on club probabilities and the flight model? \\
Placement and triggering & US7641565B2; US10247553B2 &
  Does detection reliably capture the interval needed for the requested output? \\
Staged display & US9242158B2 &
  Which values are provisional, and when does the displayed trajectory change? \\
\bottomrule
\end{tabular}
\end{table}

\begin{keypoint}
A launch monitor is a chain of observations and inferences. Radial speed,
image displacement, surface rotation, launch conditions and predicted
flight occupy different places in that chain. Combining sensors can add
information, but only after timing, coordinates and object identity agree;
it does not turn repeated use of the same data into independent evidence.

For example, Topgolf's \patent{US12544624B2} uses trajectory observations, spin candidates and probabilities associated with a likely club.
That approach differs from tracking rotating surface features.
A plausible spin value may partly reflect the prior model, so evaluating both methods requires experiments that expose their different ambiguities.

When using these outputs to investigate a swing, retain the measurement
frame, time interval, calibration and model assumptions. Launch conditions
constrain the ball's initial state; they alone do not identify the golfer's
joint torques, force history or unique movement strategy. Technical rigor
means carrying those information limits through the argument, rather
than letting a precise display imply a uniquely determined cause.
\end{keypoint}
